\section{The four rank-two types}\label{sec:rank-two}

In this section \(R\) has rank two, so \(a=12\), \(\mu=6\), \(b=12\),
and \(C\) is one of the 56 classes of \((9,2^6,\geq6)_4\) codes.

\subsection{Weight profiles and counting}

For the types with local weights \((6,6,8)\) and \((7,7,8)\), the local
plane contains two words of weight six or seven. Equality therefore
holds in \eqref{eq:light-budget}.
Every stabilizer word outside \(R\) has weight at least eight.
For type \((6,8,8)\), the possible global profiles are
\((A_6,A_7)=(1,0),(1,1),(2,0)\).
After removing the known local weight-six word, the allowed counts outside \(R\) are \((0,0),(0,1),(1,0)\).
These restrictions count distinct stabilizer words.

The adapted \(Z\)-only quotient in each rank-two type has dimension six.
Every nonzero class has minimum weight at most four, since addition of \(Z^8\)
complements a \(Z\)-only word's support. The extension map \(M\) must be
injective on this subspace at distance six. Classes of minimum weight one
must be sent to labels of coset distance at least five by
\eqref{eq:exact}, and injectivity makes these labels distinct.
This gives the following necessary counting conditions: type
\((6,6,8)\) requires four distinct labels of coset distance at least
five, and the next two types require two. Of the 56 shortened-code
classes, these tests leave 48, 52 and 52, respectively.

\subsection{The three mixed types}

For the types \((6,8,8)\) and \((7,7,8)\), the second generator
\(\zeta\) has letter \(X\) on six sites, and permutations of those six
sites fix both generators of \(R\). The searches for these two types
use this six-site action. The fixed first image is unchanged and the
other images are differences from an anchor among six points. Physical
permutations reorder and reanchor these points. A lexicographically
least full tuple is increasing and no translated anchor tuple is
smaller. The search assigns points in increasing order, so its current
tuple is a fixed initial segment of every completion of that branch.
Under a change of anchor, adding future points can only decrease the
initial entries of the sorted translated list. If the already assigned
points give a smaller initial segment after translation, every
completion has a physically equivalent tuple smaller than itself. Such a
branch cannot contain the least representative and can be discarded.
The same physical action is transported to the complementary basis,
including its possible shear. Optional relaxed complementary-domain
filters that are not invariant under reanchoring are disabled in that
search; the complete \(Y\)-test is retained.

\begin{proposition}[Mixed rank-two types]\label{prop:17-mixed}
The three local types with weights \((6,6,8)\), \((6,8,8)\), and
\((7,7,8)\) admit no extension to the hypothetical code.
\end{proposition}
\begin{proof}[Computer-assisted proof]
For each remaining shortened code and allowed pair \((A_6,A_7)\),
enumerate every
eligible first syndrome image and every compatible independent tuple of
six \(Z\)-only images. For each tuple, solve every complementary-image
assignment, including all diagonal freedom, using
\(\Gamma=MJM^\tr\) and all distance inequalities in the original coordinates.
Completed maps must also have the prescribed numbers of weight-six and
weight-seven words outside \(R\). The equations and pruning conditions
are those of Sections~\ref{sec:completion}--\ref{sec:procedure}.

The complete records contain 641 first-image cases for type \((6,6,8)\),
650 for \((7,7,8)\), and 1,950 for the three profiles of \((6,8,8)\).
Every case is exhausted with no accepted completion. For the first type,
no complete assignment to the six \(Z\)-only images passes the
necessary conditions. The six-site permutation reduction used in the
latter two searches preserves every complete extension, by the argument
above. The counting conditions and the completed searches together cover
all 56 shortened-code classes in each applicable profile.
\end{proof}

\subsection{The full plane}

For the full plane \(R=\langle Z^8,X^8\rangle\), the local
centralizer has additional structure. For a word \(v\) on \(T\), let
\(n_I,n_X,n_Y,n_Z\) be its four symbol counts. Adding \(X^8\),
\(Y^8\), or \(Z^8\) permutes these counts, so the four words of
\(v+R\) have weights \(8-n_I,8-n_X,8-n_Y,8-n_Z\), with total 24.
Every coset outside \(R\) in the local centralizer
\(H=S^\perp\cap V_T\) lies in \(S^\perp\setminus S\), so its four
words have weight at least six. All four weights are therefore six, and
every symbol count is two. The full local centralizer \(H\) is even, and
hence isotropic by \eqref{eq:polar}.
Its quotient \(K=H/R\) is an isotropic logical subspace, so \(q\leq3\)
by Proposition~\ref{prop:kernel}. The identities of that proposition,
with \(t=0\), give
\[
 h=2q,\qquad \dim N=q.
\]
The dimension-six shortened-code classification has hull distribution
\(h=0,2,6\) with 15, 33 and 8 classes; no hull-four class occurs.
Only \(q=0,1,3\) remain.

The possible local spaces \(H\) are classified as follows.
For \(q=1\), a word of \(H\setminus R\) has two occurrences of
each Pauli letter; site permutations give one local type \(H\).
For \(q=3\), every local \(H\) contains a dimension-four subcode
containing \(R\). The complete eight-site distance-six dimension-four
classification of Section~\ref{sec:local-classifications} contains three
isotropic classes with a full-support plane.
Testing all 256 normalizer cosets for each gives one dimension-five
local \(H\) type containing \(R\). In both cases \(R\) consists of
the three nonzero full-support words, so physical equivalences preserve
the plane as a whole. The search for this type uses the corresponding
physical eight-point action, verified on the full quotient basis.

\begin{proposition}[Full rank-two plane]\label{prop:17-plane}
The local type \(R=\langle Z^8,X^8\rangle\) admits no extension.
\end{proposition}
\begin{proof}[Computer-assisted proof]
For every shortened-code class of nonzero hull, enumerate every
\(q\)-space \(N\leq\rad(P)\), up to verified physical suffix automorphisms.
A space with more than two words of weight six or seven violates
\eqref{eq:light-budget} and is excluded directly. The remaining cases
have 63 \(N\)-orbits for \(q=1\), and 164 for \(q=3\).
All identifications between the quotient radical and the local kernel
radical are considered, modulo physical prefix equivalence. For \(q=3\),
the induced physical group on \(K\) has order 21. Its action on
the 168 identifications is free: an invertible identification fixed by
an element forces that element to be the identity. There are therefore
eight orbits, each of size 21. There are \(63+164\cdot8=1375\) extension problems with fixed
kernels and radical identifications, covering all 41 nonzero-hull
shortened-code classes. The \(N\)-orbits are obtained by enumerating all three lines
in each hull-two space and all 1,395 three-dimensional subspaces in each
hull-six space. The recorded orbits form a disjoint partition of these
sets; their sizes sum to three or 1,395 for each shortened code.
Each problem is searched as in Section~\ref{sec:search}, with the
prescribed kernels of Lemma~\ref{lem:fixedkernel}.
The search exhausts all assignments to the \(Z\)-only basis images.
For each of the 2,960 surviving tuples, exhaustive search finds no
assignment to the complementary images.

For the 15 hull-zero shortened-code classes, \(q=0\), and the
verification considers every admissible pair \((A_6,A_7)\), including
\(A_6=A_7=0\) and the cases with larger minimum stabilizer weight. These
are covered by separate profile exclusions or by completed extension
searches with 72,827 first-image cases. The verification program checks
both sets of records and their coverage of the allowed pairs. Together
these cover all 56 shortened-code classes and all allowed global profiles.
\end{proof}

Propositions~\ref{prop:17-mixed} and \ref{prop:17-plane} exclude every
rank-two type. Since the choice of the weight-eight word \(p\) was
arbitrary, we obtain:

\begin{corollary}\label{cor:rank-one-everywhere}
In a hypothetical \(\qc{17,3,6}\) code, every weight-eight stabilizer
word \(p\) satisfies \(S\cap V_{\operatorname{supp}p}=\langle p\rangle\).
\end{corollary}
